\section{Métodos}
\label{sec:metodos}

\subsection{Flujo general}

Todo el procesamiento es clásico y determinista: ningún paso de la segmentación
ajusta parámetros a partir de etiquetas. Cada imagen recorre la misma cadena:

\begin{enumerate}
    \item \textbf{Lectura} sin alterar valores; reducción a un canal verificada.
    \item \textbf{Reducción} a un lado de trabajo de 512 píxeles.
    \item \textbf{Campo irradiado}: región con señal del detector.
    \item \textbf{Realce} (ninguno, ecualización global o CLAHE).
    \item \textbf{Umbralización} (fija, Otsu o adaptativa).
    \item \textbf{Morfología fija}: cierre, relleno de huecos, apertura y selección
    de las dos componentes mayores.
    \item \textbf{Descriptores} sobre la máscara obtenida y sobre la de referencia.
    \item \textbf{Comparación}: Dice, Jaccard y error de cada descriptor.
\end{enumerate}

Todos los análisis respetan la partición congelada por paciente. Los que
eligen algo (estrategia, predictores, modelo) lo hacen con entrenamiento o
validación, nunca con la prueba.

\subsection{Región de interés: el campo irradiado}

Montgomery entrega siempre 4020 por 4892 píxeles con independencia de la
colimación, y el área no irradiada queda como negro puro. Dividir el área pulmonar
por el área de la imagen mediría la apertura del colimador, no la anatomía. Un
recorte rectangular no basta: deja entre el 20 y el 32\,\% de fondo, porque el
área irradiada no es rectangular.

El campo irradiado se obtiene por umbral sobre el nivel de fondo (8) y selección de
la componente conexa mayor, \textbf{sin relleno de huecos ni cierre}: las axilas
atenúan tan poco que se registran al nivel del fondo, y rellenarlas haría que el
denominador midiera el colimador y no el cuerpo. Sobre esta región se calculan el
histograma, el umbral de segmentación y el denominador de la fracción de área.


\subsection{Estrategias de segmentación (OE3)}
\label{sec:met-oe3}

Se comparan nueve combinaciones de tres realces y tres umbralizaciones. La
morfología es la misma en todas, para que una diferencia de resultado pueda
atribuirse al realce o al umbral.

\begin{table}[H]
\centering
\small
\caption{Componentes de las nueve estrategias y sus parámetros}
\begin{tabular}{@{}p{3.4cm}p{3.6cm}p{6.6cm}@{}}
\toprule
\textbf{Etapa} & \textbf{Opción} & \textbf{Parámetros} \\
\midrule
Realce & \texttt{base} & Sin realce (línea base) \\
       & \texttt{eq} & Ecualización global por función acumulada \\
       & \texttt{clahe} & CLAHE, límite de recorte 0,02, rejilla de 8 \\
\addlinespace
Umbral & \texttt{fijo} & Percentil 35 de la intensidad del campo \\
       & \texttt{otsu} & Umbral de Otsu sobre el campo \\
       & \texttt{adap} & Adaptativo local, bloque de 0,25 del lado \\
\addlinespace
Morfología (fija) & --- & Cierre (0,020 del lado), relleno de huecos, apertura
(0,012 del lado) y dos componentes mayores \\
\bottomrule
\end{tabular}
\end{table}

Una estrategia se nombra \texttt{realce\_umbral}; por ejemplo, \texttt{base\_adap}
es sin realce y umbral adaptativo. La concordancia con la referencia se mide con:
\begin{equation}
\mathrm{Dice}(A,B) = \frac{2\,|A \cap B|}{|A| + |B|}, \qquad
\mathrm{Jaccard}(A,B) = \frac{|A \cap B|}{|A \cup B|}.
\end{equation}
Una segmentación es \emph{elegible} si la imagen fue aceptada y la máscara tiene
exactamente dos componentes, una por pulmón.

La estrategia se eligió en entrenamiento por Dice mediano y consistencia entre
orígenes. En esta fase se verificó en validación: Dice por imagen de las nueve
estrategias en las 110 imágenes con máscara, ranking por origen, y comparación
pareada de la elegida frente a cada otra con la prueba de rangos con signo de
Wilcoxon y corrección de Holm sobre las 16 comparaciones.

\subsection{Descriptores}
\label{sec:met-descriptores}

Los cinco descriptores son adimensionales, porque Shenzhen no documenta el
espaciado de píxel. Se calculan igual sobre la máscara obtenida y sobre la de
referencia, y ninguna función recibe la etiqueta clínica.

\begin{table}[H]
\centering
\small
\caption{Definición de los descriptores}
\begin{tabular}{@{}p{3.5cm}p{4.6cm}p{5.6cm}@{}}
\toprule
\textbf{Descriptor} & \textbf{Fórmula} & \textbf{Qué mide} \\
\midrule
Compacidad & $C = \dfrac{4\pi A}{P^2}$ & Regularidad del contorno: menor cuanto
más irregular \\
\addlinespace
Relación de aspecto & $R = \dfrac{\text{alto}}{\text{ancho}}$ de la caja envolvente &
Proporción vertical frente a horizontal del campo pulmonar \\
\addlinespace
Fracción de área & $F = \dfrac{A_{\text{pulmón}}}{A_{\text{campo}}}$ & Proporción
del campo irradiado ocupada por pulmón \\
\addlinespace
Simetría izq-der & $S = \dfrac{|A_i - A_d|}{A_i + A_d}$ & Diferencia de área entre
los dos campos; cero si son iguales \\
\addlinespace
Energía wavelet & $E = \dfrac{1}{n}\sum_{R} \left(c_H^2 + c_V^2 + c_D^2\right)$ &
Energía media de los detalles de primer nivel (db2) dentro de la región \\
\bottomrule
\end{tabular}
\end{table}

Detalles de cálculo que condicionan la interpretación:

\begin{itemize}
    \item \textbf{Perímetro.} $P$ se cuenta como los píxeles que la región pierde al
    erosionarla con vecindad 4. Sobre una rejilla discreta, un disco da $C \approx
    1{,}24$ y no 1. El sesgo es el mismo para la máscara obtenida y la de
    referencia, de modo que su comparación es interpretable. Un borde dentado
    tiene muchos más píxeles de borde, y por eso reduce la compacidad.
    \item \textbf{Campos izquierdo y derecho.} Si la máscara tiene dos componentes,
    se ordenan por la posición horizontal de su centroide; si tiene una sola, se
    parte por la mitad de su extensión horizontal. Por convención radiológica, el
    pulmón izquierdo del paciente queda a la derecha de la imagen.
    \item \textbf{Energía wavelet.} La imagen se normaliza a $[0,1]$, se anula fuera
    de la región y se descompone con la wavelet db2 en un nivel. La energía de
    los tres detalles se promedia sobre los $n$ coeficientes dentro de la región
    reducida. Anular el exterior crea un borde artificial en el contorno, que
    aporta energía de detalle: el descriptor refleja en parte el contorno de la
    máscara y no solo la textura del parénquima.
\end{itemize}

El error de cada descriptor frente a su referencia se expresa como error relativo
con signo:
\begin{equation}
e = \frac{d_{\text{obtenido}} - d_{\text{referencia}}}{|d_{\text{referencia}}|}.
\end{equation}
Un valor negativo indica que la segmentación automática subestima el descriptor.

\subsection{Realce, ruido, frecuencia y wavelets (OE2)}
\label{sec:met-oe2}

OE2 se evaluó sobre una \textbf{muestra controlada} de 40 imágenes de
entrenamiento: 10 por cada combinación de origen y clase, con semilla fija. Cada
estrategia se mide con tres métricas dentro del campo irradiado:

\begin{itemize}
    \item \textbf{Contraste RMS}: desviación estándar de la intensidad.
    \item \textbf{Ruido estimado}: desviación del ruido estimada por la mediana
    absoluta de los coeficientes wavelet de escala fina.
    \item \textbf{Conservación de bordes}: correlación entre las magnitudes de Sobel
    de la imagen original y la procesada.
\end{itemize}

El criterio de la formulación (conservar una estrategia solo si mejora al menos una
métrica sin degradar las otras) se volvió operativo con umbrales fijados antes de
ver resultados: el contraste se degrada si baja más del 5\,\%, el ruido si sube más
del 5\,\%, y los bordes si su conservación cae por debajo de 0,90.

\paragraph{Cambios respecto a la formulación.} La formulación definía el ruido
como la desviación estándar del fondo, pero el fondo de Montgomery es relleno negro
puro y daría cero; se usó el estimador por wavelets. La conservación de bordes se
definió como correlación de gradientes en todo el campo, y no como razón sobre el
contorno de referencia, para no depender de la máscara manual.

Las estrategias se agruparon en cinco lotes:

\begin{enumerate}
    \item \textbf{Espacial}: media 5$\times$5, mediana 5$\times$5, gaussiano
    ($\sigma = 1{,}5$) y \emph{unsharp masking}; operadores de Sobel, Prewitt,
    Roberts, Laplaciano y \emph{emboss}.
    \item \textbf{Frecuencia}: pasa-bajas gaussiano y Butterworth (corte 0,15
    ciclos/píxel, orden 2), pasa-altas (0,05) y pasa-banda (0,03 a 0,15). La
    firma de cada origen se midió como la fracción de potencia por encima de 0,25
    ciclos/píxel, sobre un recorte central de 192 píxeles dentro del campo con
    ventana de Hann, para que el relleno negro no contaminara el espectro.
    \item \textbf{Wavelets}: descomposición db2 de tres niveles sobre el mismo
    recorte, con la energía relativa de cada subbanda (A3 y H, V, D de cada nivel),
    y transformada continua (sombrero mexicano) sobre un perfil horizontal.
    \item \textbf{Intensidad}: especificación del histograma de Montgomery hacia
    una referencia de Shenzhen dentro del campo, ecualización global, CLAHE
    (límite 0,01), transformaciones puntuales y requantización.
    \item \textbf{Geometría}: conectividad 4 y 8, transformada de distancia, pasos
    de la morfología y \emph{top-hat} / \emph{black-hat}.
\end{enumerate}

Las diferencias entre orígenes se cuantificaron con el delta de Cliff y su
intervalo bootstrap, con corrección de Holm cuando hubo varias pruebas.

\subsection{Desplazamiento de los descriptores (OE4)}
\label{sec:met-oe4}

Se calculó el error con signo de cada descriptor de \texttt{base\_adap} frente a
su referencia, en las imágenes elegibles de las tres particiones, y se resumió por
descriptor, partición y origen con la mediana del error con signo (sesgo), la
mediana del error absoluto (magnitud), su intervalo intercuartílico y el porcentaje
de imágenes sobreestimadas. Para la simetría, cuyo valor de referencia puede ser
casi cero, se reporta además la diferencia absoluta, porque el error relativo
explota cuando el denominador tiende a cero.

\subsection{Diferencias entre grupos (OE5)}
\label{sec:met-oe5}

El protocolo se congeló en \texttt{config/analisis\_oe5.json} antes de ver
resultados, con los hashes de las tablas que se iban a analizar.

\begin{itemize}
    \item \textbf{Unidad}: el paciente; las imágenes de un mismo paciente se
    promedian.
    \item \textbf{Población}: pacientes elegibles de \texttt{base\_adap} en las tres
    particiones. OE5 no selecciona ningún modelo, de modo que no requiere una
    partición reservada.
    \item \textbf{Estratos}: Shenzhen y Montgomery por separado, por el efecto de
    sitio.
    \item \textbf{Prueba}: Mann-Whitney $U$ bilateral, que no supone normalidad.
    \item \textbf{Tamaño de efecto}: delta de Cliff,
    \begin{equation}
    \delta = P(X > Y) - P(X < Y) = \frac{2U}{n_1 n_2} - 1,
    \end{equation}
    entre $-1$ y $1$, positivo si el grupo con hallazgos tiende a valores mayores,
    con intervalo bootstrap percentil al 95\,\% (2000 réplicas por paciente).
    \item \textbf{Comparaciones múltiples}: Holm sobre las 10 pruebas
    (5 descriptores por 2 orígenes), $\alpha = 0{,}05$.
\end{itemize}

El análisis principal usa los descriptores de la máscara de referencia. Como
sensibilidad se repitió con los descriptores automáticos de los mismos pacientes y
con las nueve estrategias en entrenamiento. Una conclusión se considera
\emph{estable} si referencia y automática coinciden en la decisión tras Holm y en
el signo del delta.

\subsection{Modelado supervisado}
\label{sec:met-modelo}

La formulación no planeaba un modelo predictivo, porque la pregunta es descriptiva.
El laboratorio de modelado del curso lo exigió, y se formuló una pregunta
complementaria: si los descriptores automáticos permiten estimar la fracción
pulmonar de la máscara manual en pacientes no vistos.

\begin{table}[H]
\centering
\small
\caption{Contrato de modelado (\texttt{config/modelado\_lab005.json})}
\begin{tabular}{@{}p{4cm}p{9.6cm}@{}}
\toprule
Unidad & Paciente (\texttt{grupo\_id}) \\
Salida & \texttt{ref\_fraccion\_area}, de la máscara manual \\
Predictores & \texttt{fraccion\_area}, \texttt{energia\_wavelet\_d1} y
\texttt{simetria\_izq\_der}, de la máscara automática \\
Excluidos & Toda columna \texttt{ref\_} o \texttt{err\_}, Dice, Jaccard,
identificadores, origen y etiqueta clínica \\
Pacientes & \ModNTrain{} en entrenamiento, \ModNVal{} en validación y \ModNTest{} en
prueba \\
Métrica primaria & RMSE \\
Referencia sencilla & Media de entrenamiento \\
\bottomrule
\end{tabular}
\end{table}

Los predictores se eligieron solo con entrenamiento: asociación con el objetivo,
mismo signo en ambos orígenes y redundancia baja entre ellos. Con $z_j$ el
predictor estandarizado con la media y la desviación de entrenamiento, se
compararon:
\begin{align}
\text{Grado 1:}\quad \hat{y} &= \beta_0 + \sum_{j} \beta_j z_j, \\
\text{Grado 2:}\quad \hat{y} &= \beta_0 + \sum_{j} \beta_j z_j
  + \sum_{j \le k} \beta_{jk}\, z_j z_k .
\end{align}
El grado 2 es no lineal en los predictores y lineal en los coeficientes, que se
estiman por mínimos cuadrados. La imputación, el escalamiento y la expansión se
ajustan dentro de un \emph{pipeline} solo con entrenamiento. El candidato se eligió
por menor RMSE de validación, el protocolo se congeló y la prueba se extrajo y
evaluó una sola vez. La incertidumbre del error en prueba se estimó con un
bootstrap de 2000 réplicas por paciente, con el modelo fijo.

\paragraph{Práctica común.} Con datos sintéticos de mecanismo conocido se
desarrollaron las dos rutas del laboratorio: regresión de grados 1 y 2 y
regresión logística,
\begin{equation}
P(y = 1 \mid z) = \frac{1}{1 + e^{-(\beta_0 + \beta^\top z)}},
\end{equation}
ajustada minimizando la pérdida logarítmica con penalización L2 de intensidad
inversa $C$ (solucionador L-BFGS), con $C \in \{0{,}1;\ 1\}$ y umbral de decisión
0,5. La ruta binaria no aplica al proyecto: no formula detección de tuberculosis, y
la etiqueta es variable de control.